\chapter{Conclusion}\label{chap:conclusion}

Ce mémoire s'est ouvert sur une question : pouvons-nous obtenir, à partir des premiers principes, les ingrédients du spectre Raman résonant du graphène avec des défauts, de sorte que ce spectre puisse ensuite être calculé sans aucun paramètre physique ajusté ? Pour une lacune dans le graphène, la réponse est oui, à un ingrédient
près et du moins pour ce qui est de la méthode. Tous les ingrédients calculés ici sont tirés, directement ou non, de calculs DFT cohérents entre eux, menés avec la même fonctionnelle, les mêmes pseudo-potentiels et le même sous-espace de Wannier, puis ont été livrés sur les grilles denses requises par la résonance. Le maillon manquant, le couplage électron-défaut, a été construit à partir des premiers principes.

\section{Synthèse}

La contribution centrale de ce mémoire, développée au chapitre~\ref{chap:défauts}, est une méthode \textit{ab initio} qui fournit le couplage électron-défaut sur des grilles arbitrairement denses en traitant la diffusion qu'il induit à tous les ordres dans le potentiel de défaut. Les éléments de matrice $M$ sont obtenus dans une base d'ondes planes, en séparant la partie locale du potentiel de défaut de la partie non locale issue des pseudo-potentiels, et sont validés par la reconstruction des valeurs propres de Kohn-Sham à quelques dizaines de µeV près. Deux idées rendent la méthode praticable sur les grilles de résonance. Le remplissage par zéros du potentiel de défaut allège la contrainte de commensurabilité entre la super-cellule et la grille de fonctions de Wannier, révélant du même coup la localisation de $M$ dans cette base. Cette localisation permet ensuite de résoudre la matrice $T$ dans un amas de cellules autour du défaut : l'inversion, de dimension $288\,000$ sur une grille $240\times 240$, se réduit à une dimension de $145$, indépendamment de la densité de la grille.

Appliquée à une lacune non relaxée, la méthode livre d'abord une leçon d'échantillonnage : les super-cellules se séparent en deux familles selon que leur grille contient ou non les points de Dirac, et seules les premières sont exemptes d'artefacts. Elle livre ensuite trois résultats. Le couplage est dominé par sa partie locale, qui porte de $77$ à $82~\%$ de l'élément de matrice au point de Dirac ; la partie non locale, due au retrait des projecteurs de l'atome manquant, est la même pour toutes les tailles de super-cellule et les deux références du potentiel. Le taux d'amortissement électronique dû aux défauts $\Gamma^\text{ed}$ montre que la lacune est un diffuseur résonant : un état quasi-lié $\pi$ à $-0.18$~eV sous le point de Dirac, où le taux par unité de concentration culmine à $40.6$~eV contre $3.19$~eV en médiane. Enfin, l'approximation de Born surestime ce taux d'un facteur $46$ en médiane et jusqu'à $203$ à la singularité de van Hove de la bande $\pi^*$. Pour un potentiel aussi intense qu'une lacune, le traitement de la diffusion à tous les ordres par la resommation de la matrice $T$ n'est pas un raffinement théorique, mais une nécessité.

Le chapitre~\ref{chap:eph} a fourni, à l'aide du code \texttt{EPW}, le couplage électron-phonon $g$, le taux d'amortissement électronique dû aux phonons $\Gamma^\text{ep}$ et la largeur des phonons $\gamma$, avec la même structure électronique et le même sous-espace de Wannier que le chapitre~\ref{chap:défauts}. Cette cohérence permet une comparaison qu'aucune combinaison de modèles ne garantit : à $300$~K et pour une concentration de lacunes de $1~\%$, le rapport $\Gamma^\text{ed}/\Gamma^\text{ep}$ a une médiane de $0.55$ et varie de $0.13$ à $35$ sur une fenêtre de $\pm 3$~eV autour du point de Dirac, le défaut l'emportant entre $-1.49$ et $+0.74$~eV. Pour des concentrations de l'ordre du pour cent, les deux mécanismes contribuent à la largeur de résonance et aucun ne peut être négligé. Comme $\Gamma^\text{ed}$ est linéaire en concentration dans la limite diluée, les phonons dominent aux concentrations de l'ordre du dixième de pour cent, sauf au voisinage immédiat de la résonance.

\section{Limites}

Ces résultats ont toutefois des limites qui ont été documentées au fil des chapitres et sont rassemblées ici. Premièrement, la lacune n'est pas relaxée et sa reconstruction de Jahn-Teller est ignorée~\cite{el-barbaryStructureEnergeticsVacancy2003}, puis le calcul n'est pas polarisé en spin alors que la lacune porte un moment magnétique~\cite{yazyevDefectinducedMagnetismGraphene2007}. Ensuite, le taux d'amortissement dû aux défauts varie encore de $2.4~\%$ entre les super-cellules $9\times 9$ et $12\times 12$, sans plateau, et la référence du potentiel, vide ou plan, reste indéterminée. De plus, la comparaison avec Kaasbjerg~\cite{kaasbjergAtomisticMatrixTheory2020} n'est que qualitative : nos éléments de matrice intravallée sont à $10~\%$ des siens à convention égale et l'état quasi-lié est bien là où la littérature le place, sous le point de Dirac~\cite{kaasbjergAtomisticMatrixTheory2020,pereiraDisorderInducedLocalized2006}, mais le maximum de la densité d'états moyennée sur le désordre est décalé de $45$ à $60$~meV vers le point de Dirac, plus étroit et plus haut, un écart que ni nos éléments de matrice, ni le protocole d'élargissement, ni les paramètres numériques n'expliquent et que nous n'avons pas investigué. Du côté phonon, la profondeur des anomalies de Kohn dépend de l'élargissement de Marzari-Vanderbilt : à l'élargissement utilisé pour la production, elles n'amollissent les modes que de $17$ et $14$~cm$^{-1}$ à $\Gamma$ et $K$ par rapport à un calcul qui ne les échantillonne pas, sans effet au-delà du pour cent sur $g$, $\Gamma^\text{ep}$ et $\gamma$. Enfin, la fonctionnelle PBE sous-estime le couplage électron-phonon au point $K$~\cite{lazzeriImpactElectronelectronCorrelation2008}.

\section{Perspectives}

Comme tout travail de méthode, ce mémoire ouvre plus de questions qu'il n'en ferme. Nous les organisons en cercles de plus en plus larges : le spectre Raman résonant qui motive tout ce travail, les mesures qui permettraient d'éprouver la méthode sans l'attendre, les approximations du formalisme qui peuvent être levées, puis les défauts et les matériaux qu'elle peut atteindre au-delà du graphène et de la lacune.

\bloc{Vers le spectre Raman résonant}

L'étape suivante est l'assemblage des ingrédients en intensités. Elle demande d'abord d'évaluer le couplage électron-photon sur les mêmes grilles denses. Ce couplage devrait être le plus simple des trois puisque, dans l'approximation dipolaire, il ne dépend que d'un seul vecteur d'onde, ne requiert ni super-cellule ni calcul perturbatif supplémentaire, et s'obtient par interpolation par fonctions de Wannier de l'opérateur vitesse, une procédure standard~\cite{yatesSpectralFermiSurface2007}. Elle demande ensuite d'intégrer les corrections ci-dessus : corriger les bandes électroniques, le couplage électron-phonon et le couplage électron-défaut par l'approximation $GW$, relaxer la structure autour de la lacune, évaluer l'effet de la polarisation en spin sur le couplage aux énergies de résonance et pousser la convergence du taux en taille de super-cellule. Encore faut-il que la correction $GW$ suffise : sous une excitation infrarouge, le couplage électron-phonon au bord de la zone de Brillouin déduit de l'expérience dépasse nettement celui prédit par la $GW$~\cite{venanziProbingEnhancedElectronphonon2023,graziottoTheoryInfraredDoubleresonance2024}, et les effets excitoniques, absents en $GW$, renormalisent fortement les anomalies de Kohn au point $K$~\cite{guandaliniExcitonicEffectsPhonons2025}. Du côté numérique, l'intensité est dominée par les états proches de la condition de résonance. Puisque l'interpolation par fonctions de Wannier permet d'évaluer tous les ingrédients sur une grille arbitraire, des grilles non uniformes avec un échantillonnage plus dense autour de ces états et moins dense ailleurs devraient réduire considérablement le coût des sommes du chapitre~\ref{chap:Raman}.

Le premier spectre calculé devrait être celui du graphène parfait, avec ses pics G et 2D comparés à la littérature~\cite{reichardtInitioCalculation$G$2017,venezuelaTheoryDoubleresonantRaman2011}, avant d'y ajouter les défauts. Les processus phonon-défaut livreront ensuite le pic D intervallé et le pic D$'$ intravallé. Le rapport de leur intensité, qui sert expérimentalement à distinguer les types de défauts~\cite{eckmannProbingNatureDefects2012}, est une cible de choix : l'expérience le mesure et aucun calcul \textit{ab initio} ne le prédit encore. En outre, les défauts peuvent aussi diffuser les phonons. Cette diffusion se traite par une matrice $T$ construite sur la perturbation des constantes de force induite par les défauts, comme l'ont fait Polanco et Lindsay pour les phonons acoustiques du graphène~\cite{polancoInitioPhononPoint2018}. Étendre leur méthode aux phonons optiques en bord de zone donnerait leur élargissement par les défauts, qui pourrait avoir un effet sur la largeur du pic D aux concentrations de l'ordre du pour cent.

\bloc{Validations indépendantes}

Le couplage électron-défaut peut lui aussi être confronté à l'expérience sans attendre le spectre Raman : deux mesures déjà publiées portent sur des lacunes et se prêtent directement à la comparaison. La première est la mobilité électronique du graphène irradié, mesurée en fonction de la densité de lacunes~\cite{chenDefectScatteringGraphene2009}. La reproduire demande d'étendre la matrice $T$ aux propriétés de transport, et ce serait le test le plus direct de l'intensité de couplage. La seconde est la densité d'états locale au voisinage d'une lacune, mesurée par microscopie à effet tunnel à la surface du graphite~\cite{ugedaMissingAtomSource2010}. Elle découle de la fonction de Green locale et de la matrice $T$ que le schéma de l'amas du chapitre~\ref{chap:T} calcule déjà sur les orbitales voisines du défaut, et trancherait, moyennant la différence entre graphite et graphène isolé, la position de l'état quasi-lié, principal désaccord avec la littérature.

\bloc{Au-delà des approximations du formalisme}

L'extension la plus directe est la structure matricielle en bande de la matrice $T$, que nous avons ignorée et qui compte parmi les causes possibles de l'écart avec Kaasbjerg. Viendrait ensuite l'approximation de la matrice $T$ autocohérente~\cite{klauderModificationElectronEnergy1961, peresElectronicPropertiesDisordered2006}, qui inclut la rétroaction de l'élargissement de désordre sur la diffusion elle même. Négligeable dans la limite diluée, elle devient nécessaire au voisinage immédiat du point de Dirac où elle engendre une densité d'états finie. Au-delà de la limite diluée, enfin, l'approximation du potentiel cohérent~\cite{sovenCoherentPotentialModelSubstitutional1967} et des calculs explicites sur des configurations de plusieurs défauts permettraient d'identifier la concentration à partir de laquelle cette limite échoue, qui est une question ouverte en soi.

\bloc{Au-delà du graphène et de la lacune}

L'un des avantages principaux des méthodes \textit{ab initio} est qu'elles sont agnostiques au système. En effet, la méthode développée dans ce mémoire n'est liée ni au graphène, ni à la lacune. Elle se généralise sans changement de principe à d'autres défauts neutres ou subsitutions, dans n'importe quel matériaux, pourvu que le défaut tienne dans une super-cellule traitable avec la DFT et que le matériau ce prête à une interpolation par fonctions de Wannier. La méthode doit cependant être adaptée pour des défauts chargés, qui demandent un traitement spécial de leur potentiel coulombien de longue portée~\cite{luFirstprinciplesIonizedimpurityScattering2021}, puis pour des défauts magnétiques, à commencer par la lacune elle-même, qui requièrent une matrice $T$ dépendante du spin. Pour ce qui est des autres systèmes à étudier, des candidats intéressants seraient les dichalcogénures de métaux de transition, comme le MoS$_2$, puisque l'introduction de défauts y active également des pics Raman. Cependant, les excitons dominent la réponse optique de ces matériaux ; combiner la diffusion par les défauts à une théorie Raman qui inclut les effets excitoniques et non adiabatiques~\cite{reichardtTheoryResonantRaman2019,reichardtNonadiabaticExcitonphononCoupling2020} ouvrirait un champ,  à notre connaissance, encore inexploré. La méthode n'est pas non plus confinée à deux dimensions. Son extension aux matériaux tridimensionnels ne demande rien de nouveau, sinon un amas plus grand et des super-cellules plus coûteuses, un coût que le remplissage par zéros réduit à celui de contenir le potentiel du défaut.

\bigskip

Le titre de ce mémoire annonçait un chemin plutôt qu'une destination. Là où la littérature offrait deux moitiés d'une solution, l'une interpolée sur des grilles denses mais limitée au premier ordre, l'autre à tous les ordres mais sans interpolation, ce mémoire les réunit. Les ingrédients du spectre Raman résonant du graphène avec des défauts sont ainsi calculés, à un près, dans un même cadre et sans paramètre ajusté sur l'expérience. Il reste à les assembler. Le jour où ce spectre sera calculé, l'expérimentateur disposera de la référence qui manque à la mesure : une empreinte digitale du matériau, calculée défaut par défaut.